% !TeX root = ../main.tex

\xchapter{绪论}{Introduction}

\xsection{研究背景及意义}{Research Background and Significance}

强化学习（Reinforcement Learning, RL）作为解决序贯决策问题的核心机器学习范式\cite{zheng2022application}，
其基础理论建立在马尔可夫决策过程（Markov Decision Process, MDP）\cite{gao2004survey}之上，
旨在通过智能体与环境的持续交互，利用试错机制与延迟奖励信号优化行为策略，
从而实现长期累积期望回报的最大化。近年来，随着深度学习\cite{lecun2015deep}技术的迅猛发展，
深度强化学习\cite{schulman2017proximal}将深度神经网络强大的非线性特征提取与高维连续函数逼近能力引入决策框架，
在诸多具有极高复杂度的理论控制任务中取得了突破性进展。

然而，现有的大多数深度强化学习算法在理论推导与模型训练中，
普遍隐含了一个严格的假设，即训练与测试阶段的观测条件保持一致。
在理想的仿真平台中，状态转移函数与奖励函数是静态且完全可知的，
智能体所接收的观测呈现完美的连续性与无损性。但在真实的非平稳环境中，
这一假设难以成立。观测噪声、视觉特征变化以及状态维度的部分缺失等因素，
均会导致测试阶段的观测分布与训练数据发生显著偏移。当观测输入偏离训练分布时，
深度神经网络的特征表示极易失效，进而导致策略退化与价值估计发散。
这种由分布偏移引发的表征问题，在强化学习的两大分支——视觉强化学习（Visual RL）与离线强化学习（Offline RL）中，
各自衍生出了不同的理论挑战。

在高维视觉强化学习研究中，模型直接以高维图像像素作为输入，
采用端到端的学习范式提取隐层状态并输出控制动作。
这种范式虽然规避了人工特征工程的主观偏差，
但高维视觉空间中充斥着大量与底层动力学无关的冗余信息
（如光度变化、背景纹理差异、相机视角偏移等）。
卷积神经网络等视觉表征模块在联合优化时具有强烈的局部拟合倾向，
容易将这些环境中的虚假关联作为决策依据\cite{zhang2018study}。
当测试阶段遭遇视觉要素的分布偏移时，泛化能力的缺失将直接导致策略崩溃。
为了提升视觉强化学习在未见视觉场景下的鲁棒性，
学术界广泛借鉴了计算机视觉领域的数据增强（Data Augmentation）技术，
通过对训练样本施加随机平移、裁剪、色彩抖动等变换，扩充观测分布的支撑集，
迫使模型学习具备变换不变性的泛化特征。

然而，从强化学习理论来看，直接引入常规的数据增强方法会产生内在冲突：
泛化能力的提升往往以牺牲样本利用效率与训练稳定性为代价。
与监督学习中分类标签对图像变换具有绝对不变性的假设不同，
强化学习的核心优化目标是基于时间差分（Temporal Difference, TD）
方法学习精确的状态-动作价值函数。
TD自举机制对状态序列的时间连续性与马尔可夫平滑性要求极高。
当引入大尺度的数据增强操作（如剧烈的随机卷积或大角度旋转）时，
不仅会破坏输入图像中的几何拓扑信息，更会导致相邻时间步的状态表征产生非平稳跳变。
这种表征层面的高频震荡在贝尔曼更新过程中会被迅速放大，
导致 TD 目标值的估计方差急剧攀升，进而引发策略梯度更新的发散。
此外，部分增强类型在特定控制任务中具有内在的破坏性，
例如旋转操作会摧毁状态空间中的方向语义。
因此，如何在引入数据多样性的同时有效抑制高方差TD误差的传播，
要求算法具备任务感知能力，能够自适应地筛选与当前任务动力学相适配的数据增强策略。
这是当前视觉强化学习实现高效泛化亟待解决的表征优化难题。

与高维视觉输入所面临的冗余干扰挑战不同，
离线强化学习主要处理由一系列低维、结构化特征向量构成的状态空间。
离线强化学习旨在完全脱离在线交互，仅依靠预先收集的静态次优数据集推演全局最优策略。
当前主流的离线强化学习算法，如保守Q学习（CQL）与隐式Q学习（IQL），
其核心理论贡献大多聚焦于解决动作空间的外推误差问题。
这些方法通过向价值函数注入下界惩罚项，或采用非对称的期望回归技术，
有效限制了策略对未见动作的盲目查询，从而缓解了离线自举过程中的值函数高估现象。

然而，现有离线强化学习算法的保守性约束往往仅局限于动作空间，
其理论框架隐含了一个脆弱的前提：即测试阶段的观测条件与训练数据保持一致。
在实际部署中，智能体可能仅能获取原始状态空间的部分维度，
即环境从训练时的完全可观测马尔可夫决策过程（MDP）退化为带有未知掩码观测函数的部分可观测马尔可夫决策过程（POMDP）\cite{kaelbling1998planning}。
这一训练与测试之间的可观测性落差构成了一种协变量偏移\cite{kumar2020implicit}。
传统的离线强化学习模型通常采用多层感知机拟合状态到动作的映射，
缺乏对状态各维度之间内在结构关联的深度建模\cite{chen2021decision}。当部分观测维度缺失时，
由于缺乏跨维度的特征推断与信息补全能力，神经网络的内部激活值会发生剧烈偏移，
产生严重的特征坍塌现象。因此，面对部分可观测条件下的观测缺失与特征坍塌，
如何引入有效的表征正则化机制，在连续潜在空间中保留状态维度间的结构关联信息，
并通过更紧凑的训练范式实现表征学习与策略优化的高效协同，
是离线强化学习在POMDP场景下亟待深入探索的核心问题。

综上所述，无论是视觉空间中的冗余扰动，还是部分可观测条件下的观测缺失，其共性均揭示了当前强化学习算法在面对不完美观测时，表示学习能力的不足。传统的强化学习优化范式侧重于策略与价值网络的设计，而忽视了在观测输入前端构建稳健、具备推理能力的表征基础。

基于此背景，本文以提升强化学习在非平稳环境与观测分布偏移条件下的泛化能力与鲁棒性为总目标，重点围绕“面向复杂观测条件的强化学习表征优化与策略协同”展开系统性研究。在理论研究层面，本文致力于打破传统强化学习端到端盲目联合优化的理论局限，深入探索表示学习辅助机制与强化学习策略优化过程的解耦与协同机制。针对视觉强化学习的高维观测分布偏移问题，本文提出基于任务感知的数据增强自适应选择机制，揭示了增强数据分布、特征方差控制与任务奖励之间的关联。针对离线强化学习在部分可观测场景下的观测缺失问题，本文将自监督掩码建模思想引入低维连续状态空间，将掩码重构作为一种有效的表征正则化手段。

\xsection{国内外研究现状}{Research Status at Home and Abroad}

视觉强化学习要求智能体直接从高维像素观测中提取特征并做出决策。早期方法（如DQN及其变体）将卷积神经网络与强化学习目标进行端到端联合训练，但高维观测空间的复杂性与奖励信号的稀疏性限制了这类方法的样本效率。为改善这一问题，数据增强技术被引入视觉强化学习。Laskin\cite{laskin2020reinforcement}等人提出RAD框架，表明简单的随机裁剪即可提升SAC算法的数据效率。Kostrikov\cite{kostrikov2020image}等人提出DrQ算法，对同一观测施加多次随机平移增强来降低Q值估计方差，其后续版本DrQ-v2\cite{yarats2021mastering}进一步扩展至连续控制任务。Srinivas等人提出CURL\cite{srinivas2020curl}算法，将对比学习引入基于像素的强化学习，利用增强后的正样本对训练表征编码器。上述方法主要关注样本效率的提升。

随着研究的深入，视觉强化学习在未见环境下的泛化问题受到越来越多的关注。Hansen等人的SVEA\cite{hansen2021stabilizing}算法分析了强数据增强导致Q学习训练不稳定的原因，指出对增强后的观测计算TD目标会引入额外方差，并提出仅对原始观测计算目标Q值的方法。在此基础上，研究者尝试从不同角度改进增强策略。Lee\cite{lee2019network}等人提出利用随机卷积增强图像的高频成分以提升泛化性。Hansen和Wang提出SODA\cite{hansen2021generalization}方法，将随机覆盖与自监督一致性学习结合。Bertoin等人提出SGQN\cite{bertoin2022look}算法，通过显著性图引导智能体关注与任务相关的像素区域，在原始观测与增强观测之间对齐注意力分布。Yuan等人提出TLDA\cite{yuan2022tlda}方法，利用Lipschitz常数识别对任务决策敏感的像素区域，仅对任务无关区域施加增强，以保护关键信息不被破坏。Liu等人提出CG2A\cite{liu2023cg2a}方法，分析了同时使用多种数据增强时产生的梯度冲突现象，并通过梯度一致性约束来缓解不同增强之间的优化冲突。

上述工作从不同角度推进了视觉强化学习的泛化研究，但大多预设了固定的增强类型或依赖人工设计的选择规则。对于给定任务，哪些增强操作是有益的、哪些是有害的，通常需要逐任务调试。缺乏任务自适应的增强选择机制，限制了现有方法在不同任务和环境间的通用性。

在离线强化学习方面，该方向旨在仅从预先收集的静态数据集中学习策略，不与环境交互。离线设定下的核心挑战在于，数据集无法覆盖完整的状态-动作空间，导致策略在评估数据集中未出现的动作时产生外推误差。为此，一系列基于保守性约束的方法被提出。Kumar等人提出BEAR\cite{kumar2019stabilizing}算法，通过分布匹配约束限制学到的策略偏离行为策略的程度。Kumar等人提出CQL\cite{kumar2020conservative}，对数据集外的状态-动作对施加Q值惩罚，学习保守的下界价值函数。Fujimoto和Gu提出TD3+BC\cite{fujimoto2019off}，在策略梯度中加入行为克隆正则项。Kostrikov等人提出IQL\cite{kostrikov2022offline}，采用非对称期望回归拟合状态价值函数$V(s)$，完全避免对数据集外动作的查询，在多种次优数据集上取得了较好效果。此外，Chen等人提出Decision Transformer\cite{chen2021decision}，将离线强化学习重新形式化为序列建模问题，利用因果Transformer根据期望回报条件生成动作序列，为离线策略学习提供了一种新范式。

上述方法主要关注动作空间的保守性约束，其隐含前提是测试阶段的观测条件与训练数据一致。然而在实际部署中，智能体常面临部分观测维度缺失的情况，即环境从训练时的完全可观测MDP退化为具有未知掩码观测函数的POMDP，由此产生的可观测性差异会导致策略性能下降。针对这一问题，近年来出现了若干研究。Yang等人\cite{yang2022rorl}提出RORL方法，在离线强化学习框架中引入保守平滑技术。该方法对数据集邻域内的扰动状态施加策略和价值函数的平滑正则化，并结合悲观自举对这些状态进行保守价值估计，在D4RL基准的标准评估和对抗攻击实验中均取得了较好结果，但其设计主要面向连续小幅扰动，对大比例维度缺失的场景考虑较少。Yang等人\cite{yang2024dmbp}提出DMBP方法，从生成模型的角度处理观测扰动问题。该方法训练一个条件扩散模型，在推理时从受扰动或残缺的观测中恢复原始状态，并提出非马尔可夫训练目标以利用轨迹级别的时序信息来缓解逐步恢复中的误差累积。DMBP在随机噪声、对抗攻击及多维度观测缺失等多种场景下提升了离线策略的鲁棒性，但扩散模型的多步采样过程带来了额外的推理开销。

在部分可观测场景的研究方面，Gu等人\cite{gu2023order}提出ORDER框架，首次系统性地研究了离线强化学习在掩码POMDP族下的泛化问题。ORDER利用VQ-VAE为每个状态因子学习离散编码本，将连续状态映射为离散代理表征，并训练一个基于循环网络的代理编码器，从部分可观测的历史轨迹中预测完整的离散表征。但框架三阶段串行设计也带来了局限：前阶段学到的离散表征是固定的，无法根据下游决策任务的需求进行调整；此外，离散量化在状态维度间存在连续依赖关系时可能造成信息损失。如何在保持梯度隔离的同时实现更紧凑的训练流程，以及在连续表征空间中保留维度间的结构信息，仍有待进一步研究。

在表示学习与强化学习的结合方面，自监督学习为提升表征质量提供了有效手段。除上述CURL的对比学习方法外，生成式的掩码建模近年来受到广泛关注。He等人提出MAE\cite{he2022masked}（Masked Autoencoder），通过随机遮蔽图像块并重建被遮蔽部分，在视觉表征学习中取得了优异效果。这一思路逐步被引入强化学习的决策建模中。Seo等人提出MWM\cite{seo2022masked}（Masked World Models），在冻结的MAE表征之上训练世界模型，用于基于模型的强化学习。Liu等人提出MLR\cite{liu2022mask}（Mask-based Latent Reconstruction），将掩码重建作为强化学习的辅助任务。Wu等人提出MTM\cite{wu2023masked}（Masked Trajectory Models），对轨迹中的状态和动作序列进行掩码重建预训练，学到的模型可通过改变推理时的掩码模式直接用于策略生成、动力学预测等多种下游任务。Liu等人提出MaskDP\cite{liu2022maskdp}，将掩码自编码器应用于状态-动作轨迹的自监督预训练，通过重建被遮蔽的状态和动作来学习动力学信息，实现了向新任务的零样本迁移。

上述工作表明，掩码建模能够有效提取序列数据中的结构信息，是一种具有潜力的表征学习范式。然而，将自监督表征学习与强化学习的策略优化进行联合训练时，存在梯度冲突的问题：TD误差的梯度若回传至编码器，可能破坏重建任务所学到的特征结构，导致表征退化。ORDER通过三阶段串行训练实现了梯度隔离，但牺牲了表征与决策之间的动态适应性。如何在统一的训练过程中实现重建任务与策略优化的解耦，使两者既保持梯度隔离又能协同演化，是该方向的一个关键问题。

\xsection{论文主要研究内容}{Main Research Content of the Thesis}

本研究聚焦于强化学习在复杂观测环境下面临的表征脆弱与策略泛化不足问题，分别针对高维视觉观测中的冗余干扰与低维结构化观测中的部分缺失两类挑战，深入分析其对策略性能的影响机制，并提出两种具有针对性的表征优化算法，以提升强化学习在观测分布偏移条件下的鲁棒性与泛化能力，论文内容研究框架如图\ref{fig:framework}所示。两种算法分别针对不同观测条件下的关键难题进行设计，具体研究内容如下：

\begin{enumerate}[label=\arabic*）]
    \item 面向视觉强化学习的任务感知数据增强选择方法

    在视觉强化学习中，智能体需要从高维像素观测中提取状态特征并做出决策。由于图像中包含大量与任务无关的视觉冗余信息，模型在训练环境中容易过拟合特定的背景纹理和光照条件，导致在未见视觉场景下泛化性能下降。数据增强是提升泛化能力的常用手段，但不同增强操作对不同任务的影响差异较大，缺乏任务引导的增强策略在提升数据多样性的同时容易破坏状态空间中的关键信息，引入额外的TD估计方差，导致训练不稳定与样本效率下降。为解决上述问题，本研究提出一种基于任务感知的数据增强选择方法（PICK）。首先，引入一个与策略网络并行的选择网络（Pick Network），根据当前任务的状态分布动态输出各数据增强方法的选用概率，使增强策略能够随训练过程自适应调整。其次，利用Critic网络的TD损失作为反馈信号指导选择网络的参数更新，当某种增强操作导致价值估计方差增大时，其选用概率被自动降低，从而在保留数据多样性的同时抑制对当前任务有害的增强操作。实验结果表明，PICK方法在DMControlGB2基准测试中，相比现有方法在未见视觉场景下的泛化性能与训练稳定性上均有提升。

    \item 面向离线强化学习的自监督掩码重建鲁棒方法

    在离线强化学习中，策略基于完全可观测的静态数据集训练，但实际部署时智能体常面临部分观测维度缺失的情况，即环境从训练时的MDP退化为具有未知掩码观测函数的POMDP。由于现有离线强化学习模型未显式建模各状态维度之间的关联，当部分维度缺失时，输入偏离训练数据的分布，容易引发特征坍塌与策略性能的大幅下降。针对上述问题，本研究提出一种基于自监督掩码重建的鲁棒离线强化学习方法（Masked-IQL）。首先，构建鲁棒掩码表征学习机制：引入可学习掩码标记替代传统零值填充以消除零值歧义，采用随机掩码比例采样策略使编码器暴露于广泛的缺失率水平，并通过表征正则化（LayerNorm与一致性损失）在连续潜在空间中显式约束残缺输入的表征向完整输入对齐，同时引入轻量历史窗口聚合模块以有效利用时序信息处理持续缺失场景。其次，提出两阶段解耦训练框架：第一阶段通过随机掩码与重建损失、一致性损失训练编码器学习鲁棒表征，第二阶段在完整状态编码上训练IQL决策模块，并通过梯度截断机制实现两阶段间的单向信息流，确保时序差分学习的稳定收敛。实验结果表明，Masked-IQL在D4RL基准的动态缺失与因子消减等多种部分可观测场景下，相比基线方法在泛化性能和鲁棒性上取得了显著改善。
\end{enumerate}

\xsection{论文组织结构}{Thesis Organization}

本文围绕复杂观测环境下的强化学习表征优化与鲁棒性问题展开系统性研究，全文共分为五章：
\begin{itemize}
    \item 第一章 绪论：阐述研究背景及意义，梳理国内外研究现状，明确主要研究内容与组织结构。
    \item 第二章 相关理论与技术：介绍马尔可夫决策过程、SAC 算法、IQL 算法以及自监督表示学习理论。
    \item 第三章 基于任务感知的视觉强化学习数据增强选择方法：详细论述 PICK 算法的设计原理及其在视觉偏移环境下的优越性。
    \item 第四章 基于自监督掩码重建的鲁棒离线强化学习方法：阐述 Masked-IQL 算法在部分可观测条件下的鲁棒性与表征修复能力。
    \item 第五章 总结与展望：归纳创新点与贡献，并对未来研究方向进行展望。
\end{itemize}